% ===================================================================
%  BAGAN No. 8 — Panen.
%
%  Traduction, faite en 2026, en indonesien d'aujourd'hui, sur l'ido
%  de texto/io. Regles de la colonne : voir 10-bagan-01.tex. Deux
%  scenes, deux numerotations : les cles portent c1 et c2. Le
%  fac-simile ido compose « 13) » sans sa parenthese ouvrante au
%  gobelet ; t08-c2-03-1 est l'un des trois ecarts declares de
%  renvoji.py et la colonne ecrit la reference bien formee.
%
%  bunga jagung (1) EST LA MEILLEURE COINCIDENCE LEXICALE DU RELEVE.
%  L'ourdou nomme le bleuet گل گندم, « la fleur du ble » : la plante
%  par le champ ou elle pousse. L'indonesien dit bunga jagung, « la
%  fleur de mais » — exactement la meme facon de nommer, dans une
%  langue qui n'a aucun lien avec l'autre et dans un pays ou le bleuet
%  ne pousse pas. Deux langues arrivent au meme compose parce que la
%  chose se rencontre au meme endroit, et l'anglais « cornflower » y
%  arrive une troisieme fois. Ce n'est pas un emprunt : c'est la meme
%  observation faite trois fois.
%
%  LE FLEAU (40) : L'INDONESIEN FAIT LA SEPTIEME COLONNE SANS.
%  L'arabe egyptien, le marathe, le telougou, le tamoul et l'ourdou
%  ont tous du periphraser ; seul le coreen avait le mot et l'outil.
%  On bat le riz en Indonesie en le frappant contre une planche ou en
%  le foulant, jamais avec deux batons relies par une laniere : la
%  colonne ecrit tongkat perontok. Sept colonnes sans, une avec, et le
%  trou est toujours dans la CHOSE.
%
%  MAIS LES OISEAUX SONT TOUS LA, ET AVEC DES NOMS DE FOND :
%  burung kicuit la bergeronnette (19), burung branjangan l'alouette
%  (41) — l'alouette de Java est une espece a elle —, burung
%  layang-layang l'hirondelle (56), burung berkik la becasse (61), et
%  surtout burung gereja le moineau (16), « l'oiseau d'eglise ». Ce
%  dernier est un calque parfait : le moineau domestique niche dans
%  les toitures, et l'indonesien l'a nomme par le batiment ou on le
%  voit le plus, exactement comme le francais dit « moineau de
%  clocher ». Le tableau le place justement devant le pressoir, pres
%  d'une eglise.
%
%  ET UNE COLLISION QU'IL FAUT SIGNALER : kenari nomme DEUX choses
%  sans rapport. Le serin (35) est burung kenari, l'oiseau des iles
%  Canaries ; le kenari tout court est le Canarium, un grand arbre a
%  noix de l'archipel. Le noyer (57) ne peut donc pas s'ecrire « pohon
%  kenari » sans qu'on croie a l'arbre indonesien : la colonne ecrit
%  pohon walnut. C'est le premier cas du releve ou un emprunt est
%  choisi non pour combler un trou mais pour EVITER un homonyme.
%
%  Et la vendange reste la vendange : anggur, tong, alat pemeras, bir,
%  sari apel. Douzieme refus de substitution, dans un pays ou le
%  raisin ne fait pas de vin.
% ===================================================================

%%K t08-noto-1 noto
\VUnotes{143.2mm}{%
\textit{(*) Karena kata ini tidak tepat : apakah ini panen rami, panen}
\textit{anggur, panen apel? Barangkali lebih baik memakai «}
\VUgras{mesono} \textit{» yang diusulkan dalam Mondo XIV, hlm. 242.}
\textit{Tetapi sampai sekarang akar baru itu belum diterima oleh}
\textit{Akademi dan masih menunggu pembahasan menurut aturan Ido yang}
\textit{lazim.}%
}
%%K t08-apar-1 apar
\VUsaut{5.0mm}
\VUtitre{15.0pt}{60}{BAGAN No. 8}
\VUsaut{3.0mm}
\VUcentre{13.2pt}{90}{\textit{Adegan Pertama.}}
\VUsaut{3.0mm}
\VUpk{10.2pt}{40}{Panen (*).}
\VUsaut{4.0mm}
\VUpk{10.2pt}{40}{Pemandangan pedesaan.}
\VUsaut{3.0mm}

%%K t08-c1-01-1 p
1. --- Dengan datangnya \VUgras{musim panas}, pemandangan pedesaan
sekali lagi berubah. Benih yang dipercayakan kepada alur telah
bertunas ; tanaman kecil yang hijau keluar dari tanah dan berbulir.
Butir gandum terbentuk, lalu masak, dan sekarang tinggal memanennya
saja.

%%K t08-c1-02-1 p
2. --- \VUgras{Bunga-bunga pedesaan} telah mekar. \VUgras{Bunga
jagung} \textsuperscript{(1)}, \VUgras{bunga popi kecil}
\textsuperscript{(2)} dan \VUgras{bunga digitalis}
\textsuperscript{(3)} memadukan pada warna seragam ladang gandum
kontras riang \VUgras{mahkota bunganya} yang segar.

%%K t08-c1-03-1 p
3. --- Pohon-pohon buah telah penuh berdaun ; buahnya telah masak.
\VUgras{Pohon ceri} \textsuperscript{(4)} yang kecil sarat dengan
\VUgras{buah ceri} \textsuperscript{(5)} yang indah, yang dipetik dan
dimakan anak-anak dengan sangat gembira.

%%K t08-c1-04-1 p
4. --- Mulai sekarang ternak boleh menghabiskan sepanjang hari di
udara terbuka.

%%K t08-c1-04-2 p
\VUgras{Anak-anak domba} \textsuperscript{(6)} telah tumbuh besar dan
menjadi \VUgras{domba-domba betina} \textsuperscript{(7)} yang indah dan
\VUgras{domba-domba jantan} \textsuperscript{(8)} yang indah, dengan
bulu wol yang lebat. Mereka merumput dengan tenang di
\VUgras{rumput} \VUgras{padang rumput} \textsuperscript{(9)} yang
berbintik-bintik \VUgras{bunga aster kecil.}

%%K t08-c1-04-3 p
Sebentar lagi \VUgras{binatang-binatang itu} akan dicukur untuk
diambil \VUgras{wolnya,} yang darinya dibuat kain pakaian kita.

%%K t08-tit-2 sub
\VUsaut{5.0mm}
\VUpk{10.2pt}{40}{Gembala.}
\VUsaut{3.4mm}

%%K t08-c1-05-1 p
5. --- \VUgras{Gembala} \textsuperscript{(10)} tampak tidak begitu
memperhatikan mereka. Yang ia risaukan hanyalah badai yang lewat di
kaki langit, dan \VUgras{penjaga domba} \textsuperscript{(13)}, yang
mendekatkan \VUgras{labu airnya} \textsuperscript{(14)} ke bibirnya,
tampak lebih tidak peduli lagi.

%%K t08-c1-06-1 p
6. --- Panasnya menindih ; sebab \VUgras{anjing} \textsuperscript{(15)}
pun datang menghilangkan hausnya. Ia menjilat air segar \VUgras{anak
sungai} \textsuperscript{(16)} dan menakuti seekor \VUgras{katak kecil}
\textsuperscript{(17)} yang malang, yang meringkuk di rumput tepian.
Katak itu melompat ke air jernih tempat \VUgras{teratai}
\textsuperscript{(18)} mekar.

%%K t08-c1-07-1 p
7. --- Seekor \VUgras{burung kicuit} \textsuperscript{(19)} mandi di
dekat \VUgras{mata air kecil} \textsuperscript{(20)} yang memancar di
antara dua batu karang, lalu bersembunyi di bawah \VUgras{semak
berduri} \textsuperscript{(21)} dan \VUgras{semak hawthorn}
\textsuperscript{(22)}.

%%K t08-tit-3 sub
\VUsaut{5.0mm}
\VUpk{10.2pt}{40}{Panen.}
\VUsaut{3.4mm}

%%K t08-c1-08-1 p
8. --- Bagi anak-anak desa, panen gandum adalah masa yang sangat
menyenangkan. Mereka \VUgras{berjungkir balik} \textsuperscript{(24)}
dengan riang — di atas \VUgras{onggokan jerami} \textsuperscript{(23)}.
Mereka membantu \VUgras{para pemanen} \textsuperscript{(26)}, dan
sementara para pemanen memotong \VUgras{ladang gandum}
\textsuperscript{(27)} dengan \VUgras{sabit-sabit besar}
\textsuperscript{(25)} mereka yang diasah tajam pada \VUgras{batu asah}
\textsuperscript{(30)}, anak-anak mengumpulkan gandum dan mengikatnya
menjadi \VUgras{berkas} \textsuperscript{(31)} atau \VUgras{ikatan
kecil} \textsuperscript{(32)}.

%%K t08-c1-09-1 p
9. --- Kadang-kadang yang paling besar di antara mereka diizinkan
memotong dengan \VUgras{sabit} \textsuperscript{(12)}
\VUgras{batang-batang keemasan} \textsuperscript{(28)} yang membawa
\VUgras{bulir}
\textsuperscript{(29)} berat penuh butir ; dan yang kecil-kecil, sambil
menjaga \VUgras{ternak} \textsuperscript{(33)} yang merumput dengan
tenang di \VUgras{padang rumput} \textsuperscript{(34)} di bawah
naungan \VUgras{pohon linden} \textsuperscript{(35)} yang besar,
menangkap \VUgras{kupu-kupu} yang indah dengan \VUgras{jaring}
\textsuperscript{(36)} mereka.

%%K t08-c1-09-2 p
Beberapa di antaranya bahkan memanjat \VUgras{gerobak}
\textsuperscript{(37)} untuk menata berkas-berkas yang diulurkan kepada
mereka dengan \VUgras{garpu rumput} \textsuperscript{(38)}.

%%K t08-c1-10-1 p
10. --- Di seluruh \VUgras{dataran} \textsuperscript{(39)} tempat
\VUgras{burung-burung branjangan} \textsuperscript{(41)} beterbangan,
kesibukan merajalela. Semua orang sibuk dengan panen. Tetapi sementara
orang miskin terus memakai perkakas lama — \VUgras{sabit besar, sabit}
dan \VUgras{tongkat perontok} \textsuperscript{(40)} —, tuan tanah yang
kaya menyuruh memotong gandum atau \VUgras{gandum hitam}
\textsuperscript{(43)} mereka dengan \VUgras{mesin pemanen}
\textsuperscript{(42)}. Kemudian, tanpa perlu membawa berkas ke lumbung,
orang membuat di tempat itu juga \VUgras{tumpukan berkas}
\textsuperscript{(44)} yang besar.

%%K t08-c1-11-1 p
11. --- Lalu didatangkan sebuah \VUgras{mesin perontok}
\textsuperscript{(47)}, yang digerakkan oleh \VUgras{mesin uap
berjalan} \textsuperscript{(45)} melalui \VUgras{sabuk penggerak}
\textsuperscript{(46)}, dan demikianlah butir gandum terpisah dari
\VUgras{jerami} tanpa meninggalkan ladang tempat ia ditaburkan.

%%K t08-tit-4 sub
\VUsaut{5.0mm}
\VUpk{10.2pt}{40}{Badai.}
\VUsaut{3.4mm}

%%K t08-c1-12-1 p
12. --- Sering \VUgras{badai hebat} \textsuperscript{(53)} terjadi pada
musim panen. Mula-mula gemuruh \VUgras{guruh} terdengar dari jauh ;
\VUgras{angin barat yang kencang} \textsuperscript{(54)} mulai bertiup
dan, sambil mengerang, membengkokkan pohon-pohon paling besar di
\VUgras{hutan kecil} \textsuperscript{(49)}. \VUgras{Awan-awan tebal}
\textsuperscript{(56)} hitam menyerbu langit ; mereka dilintasi
zigzag \VUgras{kilat} \textsuperscript{(55)} yang menyilaukan.
\VUgras{Hujan} dan kadang-kadang \VUgras{hujan es} mulai turun,
menghancurkan panen yang siap dipotong.

%%K t08-c1-13-1 p
13. --- Sering petir membakar suatu bangunan. Demikianlah tahun lalu
atap \VUgras{kincir angin} \textsuperscript{(50)} menjadi abu dan dua
\VUgras{sayapnya} \textsuperscript{(51)} hancur berkeping.

%%K t08-c1-14-1 p
14. --- Sesudah beberapa jam, ketenangan lahir kembali.
\VUgras{Pelangi} \textsuperscript{(52)} yang cemerlang muncul di kaki
langit dan alam melanjutkan kembali hidupnya yang terputus beberapa
saat. Orang-orang desa, yang berlindung di \VUgras{gubuk-gubuk}
\textsuperscript{(48)}, kembali bekerja dan dengan gagah berusaha
memperbaiki kerusakan yang baru saja mereka alami.

%%K t08-tit-5 sub
\VUsaut{6.0mm}
\VUcentre{13.2pt}{90}{\textit{Adegan Kedua.}}
\VUsaut{3.6mm}

%%K t08-tit-6 sub
\VUpk{10.2pt}{40}{Perburuan. --- Panen anggur.}
\VUsaut{1.4mm}
\VUpk{10.2pt}{40}{Penangkapan ikan.}
\VUsaut{4.0mm}

%%K t08-c2-01-1 p
1. --- Menjelang akhir \VUgras{Agustus} atau pertengahan
\VUgras{September}, menurut daerahnya, musim berburu dibuka. Baru saja
sinar matahari pertama membuat \VUgras{kadal-kadal}
\textsuperscript{(2)} keluar dari lubangnya, ketika \VUgras{pemburu}
\textsuperscript{(5)} berangkat dengan \VUgras{anjing-anjingnya}
\textsuperscript{(4)}.

%%K t08-c2-02-1 p
2. --- Ia mengenakan \VUgras{pakaian berburunya} \textsuperscript{(9)}
yang kuat dari kain tebal, memasang \VUgras{pembalut betis} kulit,
yang dengannya ia akan dapat berjalan di rumput basah tempat
\VUgras{kodok-kodok} \textsuperscript{(1)} bersembunyi ; ia mengambil
\VUgras{senapannya} \textsuperscript{(6)} dan \VUgras{tas buruannya}
\textsuperscript{(10)}, dan inilah ia berangkat.

%%K t08-c2-03-1 p
3. --- Semangat perburuan membawanya ke tengah kebun anggur tempat
panen anggur sedang berlangsung dengan meriah. Anak-anak petani
anggur, yang biasanya menjaga kambing di jalan, hari ini membiarkan
kambing itu merumput sesukanya, dan sesudah mengikat pada sebuah
patok seekor \VUgras{kambing jantan} \textsuperscript{(3)} tua yang
pasti akan memanfaatkan kebebasannya untuk melarikan diri, mereka pun
mengambil bagian dalam kesenangan itu. Yang seorang mengambil setandan
anggur dari \VUgras{keranjang panen} \textsuperscript{(12)} dan
bersenang-senang mengalirkan \VUgras{sarinya} \textsuperscript{(14)} ke
dalam sebuah \VUgras{cangkir} \textsuperscript{(13)} untuk membuat sari
buah yang belum beragi. Yang lain memakai \VUgras{mahkota daun}
\textsuperscript{(15)} yang baru saja ia anyam.

Di dekat mereka, \VUgras{burung-burung gereja} \textsuperscript{(16)}
yang tidak tahu malu datang sampai ke depan alat pemeras untuk mencuri
buah anggur.

%%K t08-c2-04-1 p
4. --- Sementara anak-anak bersenang-senang, para lelaki bekerja.
\VUgras{Para pemetik anggur} \textsuperscript{(18)} membawa buah anggur
ke ruang bawah tanah dengan \VUgras{keranjang punggung}
\textsuperscript{(17)} ; seorang \VUgras{tukang tong}
\textsuperscript{(19)} memperbaiki \VUgras{tong-tong}
\textsuperscript{(20)}, memeriksa apakah \VUgras{papan-papan tong}
\textsuperscript{(22)} dan \VUgras{dasar-dasarnya}
\textsuperscript{(21)} dalam keadaan baik, dan mengganti \VUgras{simpai
besi} \textsuperscript{(23)} yang patah. Ia pula yang mengerjakan
\VUgras{bak-bak besar} \textsuperscript{(25)} tempat orang menuangkan
\VUgras{buah anggur} \textsuperscript{(26)} supaya beragi.

%%K t08-c2-05-1 p
5. --- Setelah peragian selesai, anggur dituang dan ampasnya
diletakkan di atas \VUgras{alat pemeras} \textsuperscript{(28)} ;
dengan mengencangkan \VUgras{sekrup besar} \textsuperscript{(29)}
sedikit demi sedikit, orang memeras sarinya, yang mengalir ke dalam
\VUgras{bak kecil} \textsuperscript{(32)}, dan demikianlah diperoleh
\VUgras{anggur} \textsuperscript{(31)}.

%%K t08-tit-8 sub
\VUsaut{6.0mm}
\VUpk{10.2pt}{40}{Bir dan sari apel.}
\VUsaut{3.4mm}

%%K t08-c2-06-1 p
6. --- Pada dinding \VUgras{ruang bawah tanah} \textsuperscript{(27)}
memanjat sebatang \VUgras{hop} \textsuperscript{(30)}. Di sana ia hanya
tanaman hias ; tetapi di beberapa negeri, tempat pohon anggur sangat
jarang, hop ditanam untuk membuat \VUgras{bir.}

%%K t08-c2-07-1 p
7. --- \VUgras{Pohon apel} \textsuperscript{(33)} yang kecil sarat
dengan apel yang, bila masak, akan menghasilkan \VUgras{sari apel}
yang sangat baik. Di dalam \VUgras{sangkarnya} \textsuperscript{(34)},
\VUgras{burung kenari} \textsuperscript{(35)} dan \VUgras{burung
goldfinch} \textsuperscript{(36)} memandang dengan mata iri pada
burung-burung gereja yang bermain dengan bebas.

%%K t08-c2-08-1 p
8. --- Sejauh mata memandang, di lereng bukit, berjajar
\VUgras{tiang-tiang anggur} \textsuperscript{(40)} yang menopang
\VUgras{batang-batang anggur} \textsuperscript{(39)} yang sangat indah, sarat dengan
\VUgras{tandan} yang berat.

%%K t08-c2-08-2 p
Sepanjang tahun \VUgras{petani anggur} \textsuperscript{(42)} bekerja
merawat \VUgras{kebun anggurnya} \textsuperscript{(38)}, dan sekarang
ia diganjar jerih payahnya dengan panen yang bagus. \VUgras{Para
pemetik perempuan} \textsuperscript{(41)} mengisi bak-bak yang
diletakkan di atas gerobak yang ditarik \VUgras{bagal}
\textsuperscript{(43)}, dan semua orang bergembira.

%%K t08-tit-9 sub
\VUsaut{5.0mm}
\VUpk{10.2pt}{40}{Penangkapan ikan.}
\VUsaut{3.4mm}

%%K t08-c2-09-1 p
9. --- Demikian pula halnya dengan \VUgras{para pemancing}
\textsuperscript{(44)}, yang sejak pagi datang duduk di tepi air dengan
\VUgras{joran} \textsuperscript{(45)} mereka.

%%K t08-c2-09-2 p
\VUgras{Ikan} \textsuperscript{(47)} menyambar \VUgras{mata kail}
begitu mereka menjatuhkan \VUgras{tali pancingnya}
\textsuperscript{(46)} ke air, dan mereka hampir tidak sempat mengganti
\VUgras{umpan} sebelum korban baru menggelepar di ujung tali.

%%K t08-c2-09-3 p
Namun \VUgras{penjala} \textsuperscript{(48)}, yang dari
\VUgras{perahunya} \textsuperscript{(49)} melemparkan \VUgras{jala
tebar} ke tengah \VUgras{sungai} \textsuperscript{(50)}, menangkap jauh
lebih banyak ikan.

%%K t08-c2-10-1 p
10. --- Musim gugur adalah musim tamasya yang baik : berjalan kaki,
\VUgras{berkuda} \textsuperscript{(52)}, \VUgras{bersepeda}
\textsuperscript{(53)}, \VUgras{bermobil} \textsuperscript{(54)}.
\VUgras{Jalan-jalan} \textsuperscript{(51)} bersih dan orang
memanfaatkan hari-hari indah yang terakhir, sebab inilah sudah
\VUgras{burung-burung layang-layang} \textsuperscript{(56)} berkumpul di
atas atap untuk terbang dengan sayap terkembang penuh menuju
negeri-negeri yang lebih hangat. Dedaunan \VUgras{pohon walnut}
\textsuperscript{(57)} tua menguning, \VUgras{buah walnut} berjatuhan,
dan \VUgras{pohon-pohon} tinggi yang berkelompok seperti \VUgras{rumpun}
\textsuperscript{(58)} di atas \VUgras{bukit} \textsuperscript{(59)}
sebentar lagi akan meranggas.

%%K t08-c2-11-1 p
11. --- \VUgras{Matahari} \textsuperscript{(60)} terbit lebih lambat,
hari makin pendek, hawa dingin yang cukup menusuk mulai terasa pada
\VUgras{pagi} hari, dan inilah sudah kawanan \VUgras{burung berkik}
\textsuperscript{(61)} meninggalkan iklim kita yang tidak ramah.

\VUsaut{6.0mm}
\VUfilet{22mm}
